Um \emph{conjunto} é uma coleção de objetos, os elementos do conjunto.
Escrevemos $x \in A$ se $x$~é um elemento de~$A$. Conjuntos são
\emph{extensionais}---eles são completamente determinados por seus
elementos. Conjuntos são especificados \emph{listando} os elementos
explicitamente ou fornecendo uma propriedade que os elementos
compartilham (\emph{abstração}). Extensionalidade significa que a ordem
ou a maneira de listar ou especificar os elementos de um conjunto não
importa. Para provar que $A$ e~$B$ são o mesmo conjunto ($A = B$),
deve-se provar que todo elemento de~$X$ é um elemento de~$Y$ e
vice-versa.

Conjuntos importantes incluem os números naturais ($\Nat$), inteiros
($\Int$), racionais ($\Rat$) e reais ($\Real$), mas também
\emph{cadeias} ($X^*$) e \emph{sequências} infinitas ($X^\omega$) de
objetos. $A$~é um \emph{subconjunto} de $B$, $A \subseteq B$, se todo
elemento de $A$ é também um elemento de~$B$. A coleção de todos os
subconjuntos de um conjunto~$B$ é, ela própria, um conjunto, o
\emph{conjunto potência} $\Pow{B}$ de~$B$. Podemos formar a
\emph{união} $A \cup B$ e a \emph{interseção} $A \cap B$ de conjuntos.
Um \emph{par ordenado} $\tuple{x, y}$ consiste em dois objetos $x$
e~$y$, mas naquela ordem específica. Os pares $\tuple{x, y}$ e
$\tuple{y, x}$ são pares diferentes (a menos que $x = y$). O conjunto de
todos os pares $\tuple{x, y}$ em que $x \in A$ e $y \in B$ é chamado de
\emph{produto cartesiano}~$A \times B$ de $A$ e $B$. Escrevemos $A^2$
para $A \times A$; assim, por exemplo, $\Nat^2$ é o conjunto dos pares
de números naturais.
